\documentclass[12pt]{article}
\usepackage{amsmath, amssymb, geometry}
\geometry{margin=1in}
\setlength{\parindent}{0pt}
\setlength{\parskip}{1em}

\title{Practice Problem Set 6}
\author{ECON 101 -- Macroeconomics}
\date{Fall 2025}

\begin{document}
\maketitle

\section*{Problem 1}

The Bank of Canada is concerned that the inflation rate (3.8\%) is above the 2\% inflation-control target midpoint. 
It decides to implement \textbf{contractionary monetary policy}.



\begin{enumerate}
    \item[(a)] Explain how the Bank of Canada implements contractionary monetary policy using its \emph{key policy interest rate}. 
    Include the roles of the overnight rate target, settlement balances, and open-market operations.
    
    \item[(b)] Describe the complete \textbf{monetary transmission mechanism}, step by step, from the initial interest rate change to its effects on real GDP and inflation. 
    Your explanation should include effects on:
    \begin{itemize}
        \item consumption,
        \item investment,
        \item the exchange rate,
        \item net exports,
        \item aggregate demand.
    \end{itemize}
    
    \item[(c)] Using an AD--AS diagram, explain how contractionary monetary policy reduces inflation and brings the economy back toward the 2\% target.
\end{enumerate}

\newpage

\section*{Problem 2}

The Bank of Canada uses a simplified Taylor-rule style reaction function to guide its target overnight interest rate:

\[
i = 1 + 1.5(\pi - 2) + 0.5(Y - Y^*)
\]

where:
\begin{itemize}
    \item \( i \) = target overnight interest rate (in percent),
    \item \( \pi \) = current inflation rate (percent),
    \item \( 2 \) = inflation target midpoint (percent),
    \item \( Y - Y^* \) = output gap (percent).
\end{itemize}

Suppose new economic data show:
\[
\pi = 3.4\%, \qquad Y - Y^* = +1.5\%.
\]



\begin{enumerate}
    \item[(a)] Compute the Bank of Canada's recommended overnight interest rate using the reaction function. Show all steps.
    
    \item[(b)] State whether the resulting monetary policy stance is expansionary or contractionary, and explain why.
    
    \item[(c)] Explain how this policy rate affects:
    \begin{itemize}
        \item bank lending,
        \item long-term interest rates,
        \item aggregate demand,
        \item inflation in the medium run.
    \end{itemize}
\end{enumerate}

\hrule
\vspace{1em}

\begin{center}
\textit{End of Problems}
\end{center}

\end{document}
